\section*{Abstract}
\addcontentsline{toc}{section}{Abstract}

This thesis presents the design, implementation and characterization of a prototype pulse oximeter for non-invasive measurement of arterial blood oxygen saturation (SpO\textsubscript{2}). The system uses two multiplexed LEDs at \SI{660}{nm} and \SI{940}{nm}, a Hamamatsu S1223 photodiode, and a custom analog signal chain consisting of a transimpedance amplifier, a second-order Sallen-Key low-pass filter and a non-inverting gain stage. The optical interface is realized as a modular 3D-printed finger clip, and signal acquisition and processing are handled by an Analog Discovery 3 and a Python pipeline based on the ratio-of-ratios method.

The wavelength selection is validated through spectrometric characterization of the LEDs and through a hyperspectral setup based on a broadband Glo-bar source and a linear variable filter, which is used to estimate the relative transmission through finger tissue between \SI{800}{nm} and \SI{950}{nm}. The \SI{660}{nm} LED is measured at \SI{664.99}{nm} (FWHM \SI{18.30}{nm}) and the \SI{940}{nm} LED at \SI{951.79}{nm} (FWHM \SI{27.25}{nm}), the \SI{940}{nm} deviates from its expected value, which is relevant for the ratio-of-ratios interpretation.

In a live measurement the analog chain reproduces the expected PPG waveform on both channels and yields a heart-rate estimate of $(51.7 \pm 0.3)$~BPM and a proof-of-concept SpO\textsubscript{2} estimate of $(95.8 \pm 1.5)\,\%$. The system is not clinically calibrated against arterial blood gas analysis, and the SpO\textsubscript{2} value should therefore be interpreted as a functional demonstration rather than a clinically validated measurement. The work documents a complete signal chain from optics to digital SpO\textsubscript{2} estimate and identifies LED spectral width, wavelength-dependent scattering, optical coupling and other measurement error as the dominant sources of uncertainty.
